\exercisegroup

\textbf{1.} 证明公式

\[
\sum_ {k = q + 1} ^ {n - p + q + 1} \binom {n - k} {p - q - 1} \binom {k - 1} {q} = \binom {n} {p},
\]

其中 \(p \leqslant n\) 且 \(q < p\) （推广第 8 号命题 14 之推论中的论证）。

\begin{enumerate}
\def\labelenumi{\arabic{enumi}.}
\setcounter{enumi}{1}
\tightlist
\item
  如果 \(n \geqslant 1\) ，证明关系式
\end{enumerate}

\[
\binom {n} {0} - \binom {n} {1} + \binom {n} {2} - \dots + (- 1) ^ {n} \binom {n} {n} = 0.
\]

（在以下集合之间建立一一对应： \([1, n]\) 的具有偶数个元素的子集所构成的集合，与 \([1, n]\) 的具有奇数个元素的子集所构成的集合。区分n为偶数和n为奇数的情况。）

\textbf{3.} 证明以下关系式

\[
\begin{array}{l} \binom {n} {0} \binom {n} {p} + \binom {n} {1} \binom {n - 1} {p - 1} + \binom {n} {2} \binom {n - 2} {p - 2} + \dots + \binom {n} {p} \binom {n - p} {0} = 2 ^ {p} \binom {n} {p}, \\ \binom {n} {0} \binom {n} {p} - \binom {n} {1} \binom {n - 1} {p - 1} + \binom {n} {2} \binom {n - 2} {p - 2} - \dots + (- 1) ^ {p} \binom {n} {p} \binom {n - p} {0} = 0. \end{array}
\]

（考虑 \(p\) 元素子集（{[}1, n{]} 的子集）：它包含给定的 \(k\) 元素子集（ \(0 \leqslant k \leqslant p\) ），并使用练习2来证明第二个公式。）

\begin{enumerate}
\def\labelenumi{\arabic{enumi}.}
\setcounter{enumi}{3}
\tightlist
\item
  通过定义满足以下条件的映射 \(u\) （从 \([1, h]\) 到 \([0, n]\) 中）所成集合上的一个双射来证明第8号的命题15：
\end{enumerate}

\[
\sum_ {i = 1} ^ {h} u (x) \leqslant n
\]

到从 \([1, h]\) 到 \([1, n + h]\) 中的严格递增映射所成的集合上 .

\begin{enumerate}
\def\labelenumi{\arabic{enumi}.}
\setcounter{enumi}{4}
\tightlist
\item
  * (a) 设 E 是一个分配格，f 是 E 到交换半群 M（加法记号）的一个映射，使得
\end{enumerate}

\[
f (x) + f (y) = f (\sup (x, y)) + f (\inf (x, y))
\]

对于所有的 \(x, y\) （在 \(\mathbf{E}\) 中）。证明对每个有限子集 \(\mathbf{I}\) （\(\mathbf{E}\) 的子集）我们拥有

\[
\begin{array}{r l} f (\sup (I)) + \sum_ {2 n \leqslant \operatorname{Card} (I)} & \left(\sum_ {H \subset I, \operatorname{Card} (H) = 2 n} f (\inf (H))\right) \\ & = \sum_ {2 n + 1 \leqslant \operatorname{Card} (I)} \left(\sum_ {H \subset I, \operatorname{Card} (H) = 2 n + 1} f (\inf (H))\right) \end{array}
\]

（通过对 Card (I) 进行归纳。）*

\begin{enumerate}
\def\labelenumi{(\alph{enumi})}
\setcounter{enumi}{1}
\tightlist
\item
  特别地，设 A 是一个集合，设 \((\mathbf{B}_i)_{i\in \mathbf{I}}\) 是 A 的有限子集的一个有限族，并设 B 为诸子集 \(\mathbf{B}_i\) 的并。对于I的每个子集H，令 \(\mathbf{B}_{\mathbf{H}} = \bigcap_{i\in \mathbf{H}}\mathbf{B}_i\) . 证明
\end{enumerate}

\[
\begin{array}{r l} \operatorname{Card} (\mathrm{B}) + \sum_ {2 n \leqslant \operatorname{Card} (\mathrm{I})} & \left(\sum_ {\operatorname{Card} (\mathrm{H}) = 2 n} \operatorname{Card} \left(\mathrm{B} _ {\mathrm{H}}\right)\right) \\ & = \sum_ {2 n + 1 \leqslant \operatorname{Card} (\mathrm{I})} \left(\sum_ {\operatorname{Card} (\mathrm{H}) = 2 n + 1} \operatorname{Card} \left(\mathrm{B} _ {\mathrm{H}}\right)\right). \end{array}
\]

\begin{enumerate}
\def\labelenumi{\arabic{enumi}.}
\setcounter{enumi}{5}
\tightlist
\item
  证明公式
\end{enumerate}

\[
\begin{array}{r l} \binom {n + h} {h} = 1 + \binom {h} {1} \binom {n + h - 1} {h} - \binom {h} {2} \binom {n + h - 2} {h} + \dots \\ & + (- 1) ^ {h} \binom {h} {h} \binom {n} {h}. \end{array}
\]

（若 \(\mathbf{F}\) 表示由从 \([1, h]\) 到 \([0, n]\) 中的映射 \(u\) 组成的集合，使得

\[
\sum_ {x = 1} ^ {h} u (x) \leqslant n,
\]

对于 {[}1, h{]} 的每个子集 H ，考虑所有 \(u \in F\) （满足 \(u(x) \geqslant 1\) ，对每一个 \(x \in H\) ）组成的集合，并使用练习5。）

\begin{enumerate}
\def\labelenumi{\arabic{enumi}.}
\setcounter{enumi}{6}
\tightlist
\item
  \begin{enumerate}
  \def\labelenumii{(\alph{enumii})}
  \tightlist
  \item
    令 \(S_{n,p}\) 表示从 \([1, n]\) 到 \([1, p]\) 上的映射的个数。证明
  \end{enumerate}
\end{enumerate}

\[
\mathrm{S} _ {n, p} = p ^ {n} - \binom {p} {1} (p - 1) ^ {n} + \binom {p} {2} (p - 2) ^ {n} - \dots + (- 1) ^ {p - 1} \binom {p} {p - 1}.
\]

（注意 \(p^{n} = S_{n,p} + \binom{p}{1} S_{n,p-1} + \binom{p}{2} S_{n,p-2} + \cdots + \binom{p}{p-1}\) 并使用练习3。）

\begin{enumerate}
\def\labelenumi{(\alph{enumi})}
\setcounter{enumi}{1}
\item
  证明 \(S_{n,p} = p(S_{n-1,p} + S_{n-1,p-1})\) （采用第8号的方法，命题13）。
\item
  证明
\end{enumerate}

\[
\mathrm{S} _ {n + 1, n} = \frac {n}{2} (n + 1)! \quad \text { 且 } \quad \mathrm{S} _ {n + 2, n} = \frac {n (3 n + 1)}{2 4} (n + 2)!
\]

（考虑 \([1, n]\) 中那些逆像包含多于一个元素的元素 \(r\) ）。

\begin{enumerate}
\def\labelenumi{(\alph{enumi})}
\setcounter{enumi}{3}
\tightlist
\item
  如果 \(P_{n,p}\) 是将一个包含 n 个元素的集合划分为 p 个部分的划分数，证明 \(S_{n,p} = p!P_{n,p}\) .
\end{enumerate}

\begin{enumerate}
\def\labelenumi{\arabic{enumi}.}
\setcounter{enumi}{7}
\tightlist
\item
  设 \(p_n\) 为排列 \(u\) （一个集合 \(\mathbf{E}\) 的排列，该集合具有 \(n\) 个元素）的个数，这些排列满足 \(u(x) \neq x\) （对所有 \(x \in \mathbf{E}\) ）。证明
\end{enumerate}

\[
p _ {n} = n! - \binom {n} {1} (n - 1)! + \binom {n} {2} (n - 2)! - \dots + (- 1) ^ {n}
\]

* 从而得出 \(p_n \sim n! / e\) （当 \(n \to +\infty\) ；方法与练习7（a）相同）。

\begin{enumerate}
\def\labelenumi{\arabic{enumi}.}
\setcounter{enumi}{8}
\tightlist
\item
  \begin{enumerate}
  \def\labelenumii{(\alph{enumii})}
  \tightlist
  \item
    设E是一个具有qn个元素的集合。证明将E划分为n个每个含q个元素的子集的划分数等于
  \end{enumerate}
\end{enumerate}

\[
(q n)! / (n! (q!) ^ {n}).
\]

\begin{enumerate}
\def\labelenumi{(\alph{enumi})}
\setcounter{enumi}{1}
\tightlist
\item
  假设 \(\mathbf{E} = [1, qn]\) 。证明将 \(\mathbf{E}\) 划分为 \(n\) 个 \(q\) 元素子集（其中没有一个子集是区间）的划分数等于
\end{enumerate}

\[
\frac {(q n) !}{n ! (q !) ^ {n}} - \frac {(q n - q + 1) !}{1 ! (n - 1) ! (q !) ^ {n - 1}} + \frac {(q n - 2 q + 2) !}{2 ! (n - 2) ! (q !) ^ {n - 2}} - \dots + (- 1) ^ {n}
\]

（与习题7和8相同的方法）。

\begin{enumerate}
\def\labelenumi{\arabic{enumi}.}
\setcounter{enumi}{9}
\tightlist
\item
  令 \(q_{n,k}\) 为从 \([1, k]\) 到 \([1, n]\) 中的、满足以下条件的严格递增映射 \(u\) 的个数：对于每个偶数（相应地，奇数） \(x\) , \(u(x)\) 是偶数（相应地，奇数）。证明 \(q_{n,k} = q_{n-1, k-1} + q_{n-2, k}\) 并由此推出
\end{enumerate}

\[
q _ {n, k} = \binom {\left[ \frac {n + k}{2} \right]} {k}.
\]

\begin{enumerate}
\def\labelenumi{\arabic{enumi}.}
\setcounter{enumi}{10}
\tightlist
\item
  设E是一个有n个元素的集合，S是一个符号集合，使得S是E与一个仅含单个元素f的集合的不交并。假设f的权重为2，且E中每个元素的权重为0（第一章，附录，习题3）。
\end{enumerate}

\begin{enumerate}
\def\labelenumi{(\alph{enumi})}
\tightlist
\item
  设M是 \(L_{0}(S)\) 中的有效词集合，其中每个E中的元素恰好出现一次。证明如果 \(u_{n}\) 是M中元素的个数，则 \(u_{n+1} = (4n - 2)u_{n}\) ，并由此推出
\end{enumerate}

\[
u _ {n} = 2. 6 \dots (4 n - 6) \quad (n \geqslant 2).
\]

（这是关于非结合合成法则的n个不同项的乘积的个数）。

\begin{enumerate}
\def\labelenumi{(\alph{enumi})}
\setcounter{enumi}{1}
\tightlist
\item
  令 \(x_{i}\) 为 \(E\) 的出现在 \(M\) 的一个词中的元素中的第 \(i\) 个。证明当序列 \((x_{i})\) 给定时， \(M\) 的词的个数 \(v_{n}\) 等于 \((\frac{2n-2}{n-1})/n\) 并满足关系
\end{enumerate}

\[
v _ {n + 1} = v _ {1} v _ {n} + v _ {2} v _ {n - 1} + \dots + v _ {n - 1} v _ {2} + v _ {n} v _ {1}.
\]

\begin{enumerate}
\def\labelenumi{\arabic{enumi}.}
\setcounter{enumi}{11}
\tightlist
\item
  \begin{enumerate}
  \def\labelenumii{(\alph{enumii})}
  \tightlist
  \item
    让 \(p\) 和 \(q\) 是两个整数 \(\geqslant 1\) ，让 \(n = 2p + q\) ，让 \(E\) 为具有 \(n\) 个元素的集合，并令 \(N = \binom{n}{p} = \binom{n}{p + q}\) . 设 \((X_i)_{1 \leqslant i \leqslant N}\) （相应地， \((Y_i)_{1 \leqslant i \leqslant N}\) ）是 \(E\) 的按某种顺序排列的、具有 \(p\) （相应地， \(p + q\) ）个元素的所有子集的序列。证明存在一个双射 \(\varphi\) （{[}1, N{]} 到自身的双射），使得 \(X_{\varphi(i)} \subset Y_i\) 对于所有 \(i\) 。（方法类似于第4节练习6的方法：注意对每个 \(r \leqslant N\) ，集合 \(Y_j\) （至少包含 \(X_1, \ldots, X_r\) 之一）的个数是 \(\geqslant r\) .)
  \end{enumerate}
\end{enumerate}

\begin{enumerate}
\def\labelenumi{(\alph{enumi})}
\setcounter{enumi}{1}
\tightlist
\item
  让 \(h, k\) 是两个整数 \(\geqslant 1\) ，让 \(n\) 是一个整数，使得 \(2h + k < n\) ，让 \(\mathbf{E}\) 为具有 \(n\) 个元素的集合，并令 \((\mathbf{X}_i)_{1 \leqslant i \leqslant r}\) 是 \(\mathbf{E}\) 的不同子集的一个序列，每个子集都有 \(h\) 个元素。证明存在不同子集的一个序列 \((\mathbf{Y}_j)_{1 \leqslant j \leqslant r+1}\) （\(\mathbf{E}\) 的子集），每个子集都有 \(h + k\) 个元素，使得每个 \(\mathbf{Y}_j\) 至少包含一个 \(\mathbf{X}_i\) ，且每个 \(\mathbf{X}_i\) 包含于至少一个 \(\mathbf{Y}_j\) （通过对 \(n\) 进行归纳，利用（a）。）
\end{enumerate}

¶ 13. 设 E 是一个含有 2m 个元素的集合，设 q 是一个整数 \textless{} m，并令 F 为 P(E) 的所有满足以下性质的子集 S 的集合：若 X 和 Y 是 S 中两个不同的元素且 X ⊂ Y，则 Y --- X 至多有 2q 个元素。

\begin{enumerate}
\def\labelenumi{(\alph{enumi})}
\item
  让 \(\mathfrak{M} = (\mathrm{A}_i)_{1\leqslant i\leqslant p}\) 为 \(\mathcal{F}\) 中的元素，使得 \(p = \operatorname{Card}(\mathfrak{M})\) 尽可能大。证明 \(m - q\leqslant \operatorname{Card}(\mathrm{A}_i)\leqslant m + q\) 对于 \(1\leqslant i\leqslant p\) 。（用反证法论证。例如，假设存在指标 \(i\) 使得 \(\operatorname{Card}(\mathrm{A}_i) < m - q\) ，并考虑那些 \(\mathrm{A}_i\) （对于它们 \(\operatorname{Card}(\mathrm{A}_i)\) 具有最小可能值 \(m - q - s\) （其中 \(s\geqslant 1\) ））。设 \(\mathrm{A}_1,\ldots ,\mathrm{A}_r\) 为这些集合，令 \(\mathfrak{G}\) 为E的子集的集合，其中每个子集都是某些 \(\mathrm{A}_i(1\leqslant i\leqslant r)\) 与一个 \(2q + 1\) 个元素的、包含于 \(\mathrm{E} - \mathrm{A}_i\) 的子集的并 . 证明 \(\mathfrak{G}\) 至少包含 \(r + 1\) 个元素（参见练习12），并且如果 \(\mathrm{B}_1,\ldots ,\mathrm{B}_{r + 1}\) 是 \(r + 1\) 个不同元素（属于 \(\mathfrak{G}\) ），则其元素为 \(\mathrm{B}_j(1\leqslant j\leqslant r + 1)\) 以及 \(\mathrm{A}_i(r + 1\leqslant i\leqslant p)\) 的集合属于 \(\mathcal{F}\) ，与假设相反。）
\item
  由（a）推出元素个数 \(p\) （每个 \(\mathfrak{S} \in \mathcal{F}\) 的元素个数）满足不等式
\end{enumerate}

\[
p \leqslant \sum_ {k = 0} ^ {2 q} \binom {2 m} {m - q + k}.
\]

\begin{enumerate}
\def\labelenumi{(\alph{enumi})}
\setcounter{enumi}{2}
\tightlist
\item
  当2m或2q被替换为奇数时，建立与（a）和（b）类似的结果。
\end{enumerate}

\begin{enumerate}
\def\labelenumi{\arabic{enumi}.}
\setcounter{enumi}{13}
\tightlist
\item
  设E为具有n个元素的有限集，令 \((a_{j})_{1 \leqslant j \leqslant n}\) 为按某种顺序排列的E的元素序列，并令 \((\mathbf{A}_{i})_{1 \leqslant i \leqslant m}\) 为E的子集序列。
\end{enumerate}

\begin{enumerate}
\def\labelenumi{(\alph{enumi})}
\tightlist
\item
  对于每个指标 \(j\) ，令 \(k_{j}\) 为指标 \(i\) （满足 \(a_{j} \in \mathbf{A}_{i}\) ）的个数，并令 \(\mathbf{S}_{i} = \operatorname{Card}(\mathbf{A}_{i})\) . 证明
\end{enumerate}

\[
\sum_ {j = 1} ^ {n} k _ {j} = \sum_ {i = 1} ^ {m} s _ {i}.
\]

\begin{enumerate}
\def\labelenumi{(\alph{enumi})}
\setcounter{enumi}{1}
\item
  假设对于 E 中两个元素的每个子集 \(\{x, y\}\) ，恰好存在一个索引 \(i\) 使得 \(x\) 和 \(y\) 包含于 \(\mathbf{A}_i\) . 证明：如果 \(a_j \notin \mathbf{A}_i\) ，那么 \(\mathbf{S}_i \leqslant k_j\) .
\item
  根据(b)的假设，证明 \(m \geqslant n\) 。（令 \(k_n\) 是这些数中最小的 \(k_j\) 。证明我们可以假设，每当 \(i \leqslant k_n\) , \(j \leqslant k_n\) ，以及 \(i \neq j\) ，我们有 \(a_j \notin A_i\) 和 \(a_n \notin A_j\) 对于所有 \(j \geqslant k_n\) .)
\item
  根据(b)的假设，证明 \(m = n\) 当且仅当以下两种备选情况之一成立时：(i) \(A_1 = \{a_1, a_2, \ldots, a_{n-1}\}\) , \(A_i = \{a_{i-1}, a_n\}\) 对于 \(i = 2, \ldots, n\) ；(ii) \(n = k(k - 1) + 1\) ，每个 \(A_i\) 具有 \(k\) 个元素，且 \(E\) 的每个元素恰好属于 \(k\) 个集合 \(A_i\) .
\end{enumerate}

\begin{enumerate}
\def\labelenumi{\arabic{enumi}.}
\setcounter{enumi}{14}
\tightlist
\item
  设 E 是一个有限集，令 \(\mathfrak{L}\) 和 \(\mathfrak{G}\) 为 \(\mathfrak{P}(\mathbf{E})\) 的两个不相交的非空子集，并且让 \(\lambda, h, k, l\) 是四个整数 \(\geqslant 1\) 且满足以下性质：(i) 对每个 \(\mathbf{A} \in \mathfrak{L}\) 和每个 \(\mathbf{B} \in \mathfrak{G}\) , Card \((\mathbf{A} \cap \mathbf{B}) \geqslant \lambda\) ； (ii) 对于每个 \(\mathbf{A} \in \mathfrak{L}\) , Card \((\mathbf{A}) \geqslant h\) ； （iii） 对于每个 \(\mathbf{B} \in \mathfrak{G}\) , Card \((\mathbf{B}) \leqslant k\) ； (iv) 对每个 \(x \in \mathbf{E}\) ， \(\mathfrak{L} \cup \mathfrak{G}\) 中包含 \(x\) 的元素的个数恰好是 \(l\) . 证明 Card \((\mathbf{E}) \leqslant hk / \lambda\) 。（令 \((a_i)_{1 \leqslant i \leqslant n}\) 为E的不同元素按某种顺序排列成的序列，并且对每个 \(i\) 令 \(r_i\) 为 \(\mathfrak{L}\) 中 \(a_i\) 所属的元素的个数。证明，如果 Card \((\mathfrak{L}) = s\) 和 Card \((\mathfrak{G}) = t\) ，那么我们有
\end{enumerate}

\[
\sum_ {i = 1} ^ {n} r _ {i} \leqslant s h, \quad \sum_ {i = 1} ^ {n} (l - r _ {i}) \leqslant t k, \quad \text { 且 } \quad \sum_ {i = 1} ^ {n} r _ {i} (l - r _ {i}) \geqslant \lambda s t.)
\]

为使 Card (E) 等于 \(hk / \lambda\) ，必要且充分的条件是：对于每个 \(A \in \mathfrak{L}\) 和每个 \(B \in \mathfrak{C}\) ，我们有 Card (A) = \(h\) , Card (B) = \(k\) ,

Card \((\mathbf{A} \cap \mathbf{B}) = \lambda\) ，并且存在一个 \(r \leqslant l\) ，使得对于每一个 \(x \in \mathbf{E}\) ， \(\mathfrak{L}\) 中 \(x\) 所属的元素个数等于 \(r\) .

\begin{enumerate}
\def\labelenumi{\arabic{enumi}.}
\setcounter{enumi}{15}
\tightlist
\item
  设 E 是一个具有 n 个元素的有限集合，设 D 是 \(\mathfrak{P}(\mathrm{E})\) 的一个非空子集，并且让 \(\lambda\) ，k，l 为三个整数 \(\geqslant 1\) 且满足以下性质：(i) 若 A，B 是 D 的不同元素，则 Card \((\mathrm{A} \cap \mathrm{B}) = \lambda\) ； (ii) 对于每个 \(A \in D\) , Card \((\mathrm{A}) \leqslant k\) ； （iii） 对于每个 \(x \in E\) ， x 所属的 D 的元素个数等于 l。证明
\end{enumerate}

\[
n (\lambda - 1) \leqslant k (k - 1),
\]

并证明：如果 \(n(\lambda - 1) = k(k - 1)\) ，那么 \(\lambda = k\) 且 \(\operatorname{Card}(\mathfrak{D}) = n\) . （给定 \(a \in E\) ，令 \(\mathfrak{L}\) 为所有集合 \(A - \{a\}\) （其中 \(A \in \mathfrak{D}\) 且 \(a \in A\) ）所成的集合，并令 \(\mathfrak{C}\) 为所有 \(A \in \mathfrak{D}\) （满足 \(a \notin A\) ）所成的集合。把习题15的结果应用于 \(\mathfrak{L}\) 和 \(\mathfrak{C}\) .)

\begin{enumerate}
\def\labelenumi{\arabic{enumi}.}
\setcounter{enumi}{16}
\tightlist
\item
  让 \(i, h, k\) 为三个整数，使得 \(i \geqslant 1\) , \(h \geqslant i\) , \(k \geqslant i\) . 证明存在一个整数 \(m_i(h, k)\) 且满足以下性质：对于每个至少含有 \(m_i(h, k)\) 个元素的有限集 E，以及每个划分 \((\mathcal{X}, \mathcal{Y})\) （对集合 \(\mathfrak{F}_i(\mathbf{E})\) （E 的 \(i\) 元素子集所成的集合）的划分），不可能 E 的每个 \(h\) 个元素的子集都包含一个子集 \(\mathbf{X} \in \mathcal{X}\) ，且 E 的每个 \(k\) 个元素的子集都包含一个子集 \(\mathbf{Y} \in \mathcal{Y}\) ；换言之，若 E 的每个 \(h\) 个元素的子集包含某个 \(\mathbf{X} \in \mathcal{X}\) ，则存在 E 的一个具有 \(k\) 个元素的子集 A，使得 A 的每个 \(i\) 个元素的子集都属于 \(\mathcal{X}\) 。（用归纳法证明。证明我们可以取
\end{enumerate}

\[
m _ {1} (h, k) = h + k - 1, \quad m _ {i} (i, k) = k, \quad \text { 且 } \quad m _ {i} (h, i) = h,
\]

最后 \(m_i(h, k) = m_{i-1}(m_i(h-1, k), m_i(h, k-1)) + 1\) 。如果 \(E\) 是一个具有 \(m_i(h, k)\) 个元素的集合，如果 \(a \in E\) 并且如果 \(E' = E - \{a\}\) ，证明，若该命题为假，则 \(E'\) 的每个 \(m_i(h-1, k)\) 个元素的子集都包含一个 \(i-1\) 个元素的子集 \(X'\) ，使得 \(X' \cup \{a\} \in \mathcal{X}\) ，且 \(E'\) 的每个 \(m_i(h, k-1)\) 个元素的子集都包含一个 \(i-1\) 个元素的子集 \(Y'\) ，使得 \(Y' \cup \{a\} \in \mathcal{Y}\) .

\begin{enumerate}
\def\labelenumi{\arabic{enumi}.}
\setcounter{enumi}{17}
\tightlist
\item
  \begin{enumerate}
  \def\labelenumii{(\alph{enumii})}
  \tightlist
  \item
    设 E 是一个具有 p 个元素的有限有序集。若 m, n 是两个整数且 mn \textless{} p，证明 E 要么有一个 m 个元素的全序子集，要么有一个 n 个元素的自由子集（§ 1，习题 5）（利用 § 4，习题 5）。
  \end{enumerate}
\end{enumerate}

\begin{enumerate}
\def\labelenumi{(\alph{enumi})}
\setcounter{enumi}{1}
\tightlist
\item
  让 \(h, k\) 是两个整数 \(\geqslant 1\) 并令 \(r(h, k) = (h - 1)(k - 1) + 1\) . 设 \(I\) 是一个至少含 \(r(h, k)\) 个元素的有限全序集。证明，对于每个有限序列 \((x_i)_{i \in I}\) （由一个全序集 \(E\) 的元素组成），要么存在 \(I\) 的一个具有 \(h\) 个元素的子集 \(H\) ，使得序列 \((x_i)_{i \in H}\) 是递增的，要么存在 \(I\) 的一个具有 \(k\) 个元素的子集 \(K\) ，使得序列 \((x_i)_{i \in K}\) 是递减的。（将（a）应用于 \(I \times E\) .)
\end{enumerate}

